\def\draftversion{false}
\RequirePackage{ifthen}
\ifthenelse{\equal{\draftversion}{true}}{
  \documentclass[galley,aps,pra,10pt,amsmath,amssymb,
    superscriptaddress,nofootinbib,longbibliography]{revtex4-2}
}{
  \documentclass[twocolumn,aps,pra,10pt,amsmath,amssymb,
    superscriptaddress,nofootinbib,longbibliography]{revtex4-2}
}
\usepackage{graphicx}% Include figure files
\usepackage[usenames,dvipsnames]{color} % colors
\usepackage{bm} % bold math
\usepackage{soul} % \st    for strike-out
\usepackage{mathtools}
\usepackage{tabularx}
\usepackage{dcolumn}
\usepackage{bbold}
\usepackage{hyperref}
\usepackage{siunitx}
\hypersetup{
    colorlinks=true,       
    linkcolor=blue,          
    citecolor=blue,       
    filecolor=magenta,      
    urlcolor=blue          
}

\def\Red#1{\textcolor{red}{#1}}
\def\Blue#1{\textcolor{blue}{#1}}
\def\Green#1{\textcolor{OliveGreen}{#1}}
\def\Magenta#1{\textcolor{magenta}{#1}}

\def\scr{\scriptsize}
\ifthenelse{\equal{\draftversion}{true}}{
  \marginparwidth 2.7in
  \marginparsep 0.5in
  \newcounter{comm} 
  \def\commnext{\stepcounter{comm}}
  \def\commtext{{\bf\color{blue}[\arabic{comm}]}}
  \def\commmar{{\bf\color{blue}[\arabic{comm}]}}
  \def\dvm#1{\commnext\marginpar{\small DV\commmar: #1}\commtext}
  \def\tcm#1{\commnext\marginpar{\small TC\commmar: #1}\commtext}
  \def\tnewpage{\newpage\marginpar{\small Temporary newpage}}
  \newcommand{\seclab}[1]{\label{sec:#1}{\Red{\small\;\;[sec:~#1]}}}
  \newcommand{\eqlab}[1]{\Red{\hbox{\small\;\;[#1]}}\label{eq:#1}}
  \newcommand{\figlab}[1]{\Red{\hbox{\small\;\;[fig:~#1]}}\label{fig:#1}}
}{
  \def\dvm#1{}
  \def\tcm#1{}
  \def\tnewpage{}
  \newcommand{\eqlab}[1]{\label{eq:#1}}
  \newcommand{\seclab}[1]{\label{sec:#1}}
  \newcommand{\figlab}[1]{\label{fig:#1}}
}

\newcommand{\beq}{\begin{equation}}
\newcommand{\eeq}{\end{equation}}
\newcommand{\bea}{\begin{eqnarray}}
\newcommand{\eea}{\end{eqnarray}}
\newcommand{\nn}{\nonumber\\}
\newcommand{\eq}[1]{Eq.~(\ref{eq:#1})}
\newcommand{\Eq}[1]{Equation~(\ref{eq:#1})}
\newcommand{\eqs}[2]{Eqs.~(\ref{eq:#1}) and (\ref{eq:#2})}
\newcommand{\Eqs}[2]{Equations~(\ref{eq:#1}) and (\ref{eq:#2})}
\newcommand{\eqo}[2]{Eq.~(\ref{eq:#1}) or (\ref{eq:#2})}
\newcommand{\eqr}[2]{Eqs.~(\ref{eq:#1}-\ref{eq:#2})}
\newcommand{\Eqr}[2]{Equations~(\ref{eq:#1}-\ref{eq:#2})}
\newcommand{\fref}[1]{Fig.~\ref{fig:#1}}
\newcommand{\Fref}[1]{Figure~\ref{fig:#1}}
\newcommand{\frefs}[2]{Figs.~\ref{fig:#1} and \ref{fig:#2}}
\newcommand{\frefo}[2]{Figs.~\ref{fig:#1} or \ref{fig:#2}}
\newcommand{\sref}[1]{Sec.~\ref{sec:#1}}
\newcommand{\srefs}[2]{Secs.~\ref{sec:#1} and \ref{sec:#2}} 
\newcommand{\Sref}[1]{Section~\ref{sec:#1}}
\newcommand{\aref}[1]{Appendix~\ref{sec:#1}}
\newcommand{\ket}[1]{\vert#1\rangle}
\newcommand{\bra}[1]{\langle#1\vert}
\newcommand{\ip}[2]{\langle#1\vert#2\rangle}
\newcommand{\me}[3]{\langle#1\vert#2\vert#3\rangle}
\newcommand{\ev}[1]{\langle#1\rangle}
\newcommand{\wt}[1]{\widetilde{#1}}
\newcommand{\Tr}{\mathrm{Tr}\,}   
\newcommand{\code}[1]{\textsc{#1}}  
\renewcommand{\Re}{\mathrm{Re}\,}
\renewcommand{\Im}{\mathrm{Im}\,}

\newcommand{\parkmu}{ \partial_{k_\mu} }
\newcommand{\parmu}{ \partial_{\mu} }
\begin{document}

\title{Axion-active phonons in magnetic insulators}

\author{Trey Cole}
\affiliation{
Department of Physics \& Astronomy, Center for Materials Theory, Rutgers University, Piscataway, New Jersey 08854, USA}

\author{Daniel Seleznev}
\affiliation{
Department of Physics \& Astronomy, Center for Materials Theory, Rutgers University, Piscataway, New Jersey 08854, USA}
\affiliation{
Department of Physics, University of Texas at Austin, Austin, Texas 78712, USA}

\author{Andrea Urru}
\affiliation{
Department of Physics \& Astronomy, Center for Materials Theory, Rutgers University, Piscataway, New Jersey 08854, USA}

\author{David Vanderbilt}
\affiliation{
Department of Physics \& Astronomy, Center for Materials Theory, Rutgers University, Piscataway, New Jersey 08854, USA}

\begin{abstract}
The Chern-Simons axion angle $\theta$ is quantized to $0$ or $\pi$ in insulators possessing certain symmetries such as spatial inversion $\mathcal{P}$ or time-reversal $\mathcal{T}$.
Phonon modes that simultaneously break all symmetries protecting this quantization---which we call ``axion-active'' phonons---induce a first-order modulation of $\theta$ away from its quantized value.
We develop a first-principles methodology for identifying axion-active phonon modes through symmetry analysis and computing the induced linear response $\partial_\lambda \theta$ using a gauge-invariant four-curvature formulation evaluated on Wannierized tight-binding Hamiltonians derived from density-functional theory.
We apply this approach to two magnetic insulators with contrasting axion topology: (i) MnBi$_2$Te$_4$, an antiferromagnetic topological insulator with $\theta = \pi$ enforced by inversion $\mathcal{P}$ and the combined symmetry $\Theta = \mathcal{T} S_{1/2}$, where the axion-active phonons belong to the $T_2^-$ irreducible representation at the Brillouin-zone boundary; and (ii) Y$_2$Ir$_2$O$_7$, a pyrochlore iridate with all-in-all-out magnetic order studied in a correlation regime where $\theta = 0$, providing a non-topological counterpoint.
In Y$_2$Ir$_2$O$_7$, the axion-active phonons transform as the zone-center $\Gamma_2^-$ ($A_{2u}$) representation---a mode that is neither infrared- nor Raman-active, representing a ``silent'' modulation of the axion angle.
\end{abstract}

\maketitle

\noindent

%==========================================================================
\section{Introduction}
\seclab{intro}
%==========================================================================

The Chern-Simons axion angle $\theta$ describes an isotropic contribution to the orbital magnetoelectric coupling in three-dimensional insulators \cite{qi2008topological, essin2010orbital}.
In insulating materials that break both spatial inversion and time-reversal symmetries, a magnetic field $\mathbf{B}$ can induce an electric polarization $\mathbf{P}$, and an electric field $\mathbf{E}$ can induce a magnetization $\mathbf{M}$.
This coupling is described by a linear magnetoelectric coupling tensor $\alpha_{ij} = (\partial P_i / \partial B_j)_{\mathbf{E}=0} = (\partial M_i/\partial E_j)_{\mathbf{B}=0}$ \cite{Malashevich_2010}.
The Chern-Simons axion (CSA) coupling makes an isotropic contribution $\alpha^{\text{CS}}$ to the magnetoelectric tensor proportional to the dimensionless quantity $\theta$,
\begin{equation}
    \alpha^{\text{CS}}_{ij} = \frac{\theta e^2}{2 \pi h} \delta_{ij}.
    \eqlab{alpha-CS}
\end{equation}
The quantity $\theta$ is determined by the integral of a Chern-Simons 3-form over the Brillouin zone \cite{chern1974characteristic, Malashevich_2010, mong2010antiferromagnetic, essin2010orbital},
\begin{equation}
 \theta = -\frac{1}{4\pi} \int_{\text{BZ}} d^3k \,
\epsilon^{\mu\nu\sigma} \mathrm{Tr} \left[
    \mathcal{A}_\mu \partial_\nu \mathcal{A}_\sigma - \frac{2i}{3} \mathcal{A}_\mu \mathcal{A}_\nu \mathcal{A}_\sigma
\right],
\eqlab{CS3form}
\end{equation}
where $\mathcal{A}_\mu$ is the non-Abelian Berry connection matrix.
In general, $\theta$ is only well-defined modulo $2\pi$, reflecting a gauge ambiguity analogous to the quantum of polarization in the modern theory of electric polarization \cite{marzari2012}.
However, in the presence of symmetries that quantize $\theta$---either inversion $\mathcal{P}$, time-reversal $\mathcal{T}$, or certain composite antiunitary symmetries---the axion angle is pinned to the discrete values $\theta = 0$ or $\theta = \pi$ (mod $2\pi$) \cite{Turner2012, mong2010}.
Materials with $\theta = \pi$ are termed \emph{axion insulators}, and exhibit half-quantized surface anomalous Hall conductivity on symmetry-preserving surfaces \cite{qi2008topological, essin2010orbital}.

The electrodynamics of a homogeneous medium are unaffected by a spatially and temporally uniform $\theta$, being instead only sensitive to its spatial gradients---producing a local anomalous Hall current $\mathbf{J}(\mathbf{r}, t) \propto \mathbf{E} \times \nabla \theta(\mathbf{r}, t)$---or to its time derivative---producing a local chiral magnetic effect $\mathbf{J}(\mathbf{r}, t) \propto \mathbf{B}\, \partial_t \theta(\mathbf{r}, t)$.
These effects will be activated insofar as some parameters in the crystal Hamiltonian are spatially inhomogeneous or time-dependent, and their magnitudes will be controlled by the first derivatives of $\theta$ with respect to those parameters.
The axion field, usually treated as a static background quantity depending on the ground-state crystal Hamiltonian, has been proposed to become dynamical in certain contexts \cite{li_dynamical_2010, wang2016, ooguri2012instability, wider2020}.
Rather than possessing its own independent dynamics, $\theta$ would vary in concert with preexisting degrees of freedom of the solid such as phonons or charge-density waves.
Such situations could be especially interesting in magnetic insulators where $\theta$ is comparable to or equal to $\pi$.

There is reason to think that the computation of \emph{derivatives} of $\theta$ may be more straightforward than the computation of $\theta$ itself.
The axion angle is a 3D analog of the 1D electric polarization $\mathbf{P}$, and historically, first derivatives of $\mathbf{P}$---such as Born effective charges and dielectric constants---were computed well before the modern theory of polarization \cite{marzari2012}.
This can be understood because $\mathbf{P}$ itself is only well-defined modulo a quantum and is expressed as an integral of a gauge-dependent Berry connection, whereas its derivative $\partial_\lambda \mathbf{P}$ does not suffer from either ambiguity.
Analogously, $\partial_\lambda \theta$ can be expressed as a Brillouin-zone integral of a gauge-invariant four-curvature \cite{qi2008topological, taherinejad2015adiabatic, olsen2017surface},
\begin{equation}
\partial_\lambda \theta
= \frac{1}{16\pi}
\int_{\text{BZ}} d^3 k \,
\epsilon^{ijkl} \mathrm{Tr}\!\left(
\Omega_{ij} \Omega_{kl}
\right),
\eqlab{4curv}
\end{equation}
defined in the extended parameter space $(k_x,k_y,k_z,\lambda)$, where $\Omega_{ij}$ denotes the non-Abelian Berry curvature associated with coordinates $i,j \in \{k_x,k_y,k_z,\lambda\}$.
Unlike the Chern-Simons 3-form in \eq{CS3form}, this expression is manifestly gauge invariant and does not require the construction of a globally smooth Bloch-state gauge.

In this work, we develop and apply a first-principles methodology for computing the linear response of $\theta$ to lattice vibrations that modulate the axion angle away from its symmetry-enforced quantized values.
We refer to phonon modes that break all protecting symmetries---and thus induce a nonzero $\partial_\lambda \theta$---as \emph{axion-active phonons}.
We identify such modes through a systematic symmetry analysis using the \code{SMODES} software package, then compute the axion response from Wannierized tight-binding Hamiltonians derived from fully relativistic density-functional theory calculations.
We apply this approach to two magnetic insulators with contrasting axion topology: MnBi$_2$Te$_4$ (MBT), an antiferromagnetic topological insulator with $\theta = \pi$, and the pyrochlore iridate Y$_2$Ir$_2$O$_7$ (YIO), studied in a non-topological regime where $\theta = 0$.
This pairing allows us to contrast the axion-phonon response in a topological system, where $\theta$ departs from $\pi$, with that in a trivial system, where $\theta$ departs from $0$.

The remainder of this paper is organized as follows.
In \sref{theory} we review the theoretical framework underlying the axion angle and its derivatives.
In \sref{symmetry} we discuss the general symmetry requirements for axion-active phonons.
\Sref{MBT} and \sref{YIO} present the symmetry analysis, computational details, and results for MnBi$_2$Te$_4$ and Y$_2$Ir$_2$O$_7$, respectively.
In \sref{methods} we describe the shared first-principles methodology.
Finally, in \sref{conclusions} we summarize our findings and discuss their implications.

%==========================================================================
\section{Theoretical framework}
\seclab{theory}
%==========================================================================

\subsection{The Chern-Simons axion angle}

The Chern-Simons axion coupling arises from the topological part of the electromagnetic response of a three-dimensional insulator.
As formulated by Qi, Hughes, and Zhang \cite{qi2008topological}, the effective action contains a term
\begin{equation}
S_\theta = \frac{\theta e^2}{4\pi^2 \hbar c} \int d^3r\, dt\, \mathbf{E} \cdot \mathbf{B},
\eqlab{action}
\end{equation}
where $\theta$ is the axion angle given by \eq{CS3form}.
Variation of this action with respect to the electromagnetic potentials yields the magnetoelectric response in \eq{alpha-CS}.

The axion angle $\theta$ is a bulk quantity, well-defined for a fully gapped insulator.
Because the Chern-Simons 3-form in \eq{CS3form} is gauge-dependent, $\theta$ is only well-defined modulo $2\pi$: a change of gauge (a unitary rotation of the occupied Bloch states) can shift $\theta$ by an integer multiple of $2\pi$, corresponding to the addition of an integer Chern number to the surface Hall conductivity.
This situation is analogous to the electric polarization, which is defined modulo a quantum of polarization.

\subsection{Symmetry constraints on the axion angle}

Certain symmetries constrain $\theta$ to the values $0$ or $\pi$ (mod $2\pi$) \cite{Turner2012, mong2010}.
Specifically, any symmetry that reverses the spatial orientation or reverses time constrains $\theta = -\theta$ (mod $2\pi$), which forces $\theta = 0$ or $\pi$.
This includes:
(i) spatial inversion $\mathcal{P}$;
(ii) time-reversal $\mathcal{T}$;
(iii) any improper spatial rotation (mirror, glide, or roto-inversion); and
(iv) composite antiunitary symmetries such as $\mathcal{T}$ followed by a fractional lattice translation.
In magnetic insulators where $\mathcal{T}$ is broken by spin order, the quantization of $\theta$ can still be enforced by a combination of inversion $\mathcal{P}$ with residual magnetic symmetries.

\subsection{Derivative of the axion angle: four-curvature formulation}

To compute the linear response of $\theta$ to a perturbation parametrized by a scalar amplitude $\lambda$ (such as a phonon displacement), we work in the extended parameter space $(k_x,k_y,k_z,\lambda)$ and evaluate \eq{4curv}.
The Berry curvature in this four-dimensional space is defined as
\begin{equation}
\Omega_{ij} = \partial_i \mathcal{A}_j - \partial_j \mathcal{A}_i - i[\mathcal{A}_i, \mathcal{A}_j],
\eqlab{curvature}
\end{equation}
where $\mathcal{A}_{i,mn} = \langle u_{m\mathbf{k}} | \partial_i | u_{n\mathbf{k}} \rangle$ is the non-Abelian Berry connection and the indices $i,j$ run over $(k_x, k_y, k_z, \lambda)$.
The trace in \eq{4curv} runs over occupied bands.

In practice, derivatives with respect to the crystal-momentum coordinates $k_x, k_y, k_z$ are evaluated analytically using velocity operators obtained from the tight-binding Hamiltonian, while derivatives with respect to $\lambda$ are computed using finite differences between Hamiltonians evaluated at neighboring displacement amplitudes.
Because the integrand in \eq{4curv} involves only gauge-covariant quantities (Berry curvatures), the result is gauge-invariant and avoids the numerical instabilities associated with direct evaluation of the Chern-Simons 3-form.

%==========================================================================
\section{Axion-active phonons: general symmetry analysis}
\seclab{symmetry}
%==========================================================================

To induce a first-order change in $\theta$ away from its quantized value, a phonon mode must break \emph{all} symmetries that enforce the quantization $\theta = 0$ or $\pi$.
We refer to such modes as \emph{axion-active phonons}.
More precisely, given the set of crystallographic symmetries $\{g_i\}$ that individually constrain $\theta$ to a quantized value, an axion-active phonon mode is one whose displacement pattern breaks every $g_i$ simultaneously.
If even one protecting symmetry survives the distortion, $\theta$ remains quantized and the first-order response vanishes.

The identification of axion-active modes proceeds in several steps.
First, we determine the protecting symmetries of the undistorted magnetic crystal.
Second, we enumerate the irreducible representations of phonon modes at the relevant wavevectors (which may be at the zone center, zone boundary, or other high-symmetry points, depending on the material and the nature of the protecting symmetries).
Third, for each irreducible representation, we determine which of the protecting symmetries are broken by the corresponding displacement pattern.
Those representations that break all protecting symmetries are the axion-active irreps, and all modes transforming under such irreps are axion-active phonons.

In the two materials studied in this work, the protecting symmetries, the quantized values of $\theta$, and the resulting axion-active irreps all differ qualitatively, illustrating the richness of axion-phonon physics across different classes of magnetic insulators.

%==========================================================================
\section{MnBi$_2$Te$_4$}
\seclab{MBT}
%==========================================================================

\subsection{Crystal and magnetic structure}

MnBi$_2$Te$_4$ (MBT) crystallizes in a rhombohedral structure with space group $R\bar{3}m$ (No.~166).
The structure consists of septuple layers Te-Bi-Te-Mn-Te-Bi-Te stacked along the crystallographic $c$ axis.
The Mn ions adopt a $2+$ valence with spin $S = 5/2$ and a large magnetic moment of approximately $5\,\mu_\text{B}$.
Below the N\'eel temperature $T_N = 24$~K, the Mn$^{2+}$ moments order ferromagnetically within the $ab$ plane but antiferromagnetically along the $c$ axis in an A-type antiferromagnetic (AFM) arrangement \cite{Yan2019}.

In the magnetically ordered state, time-reversal symmetry $\mathcal{T}$ is broken.
However, the system retains spatial inversion $\mathcal{P}$ (centered at the Mn site) as well as the combined antiunitary symmetry $\Theta = \mathcal{T} S_{1/2}$, consisting of time reversal followed by a half-lattice translation along the $c$ axis \cite{mong2010}.
This half-translation maps Mn$^{2+}$ ions with spin-up to those with spin-down, so that $\Theta$ is preserved despite the magnetic order.
Together, $\mathcal{P}$ and $\Theta$ enforce the quantization $\theta = \pi$, making MBT both an antiferromagnetic topological insulator and an axion insulator \cite{Turner2012}.

At the surface, time-reversal symmetry is broken by the AFM ordering, leading to gapped surface Dirac modes and a half-quantized surface anomalous Hall conductivity.

\subsection{Symmetry analysis and identification of axion-active modes}

To induce a first-order change in $\theta$, a phonon mode must break both $\mathcal{P}$ and $\Theta$.
Because $\Theta$ involves a half-translation along $c$, phonons at finite crystal momentum along $k_z$ are natural candidates for breaking this symmetry.
We therefore focus on phonons at the $Z$ point of the hexagonal Brillouin zone of the non-magnetic crystal structure, corresponding to wavevector $(0,0,\pi/c)$.

Using the \code{SMODES} software package, we enumerate the irreducible representations of the $Z$-point little group and identify displacement patterns that break the relevant symmetries.
Among all irreps at $Z$, the $T_2^-$ representation simultaneously breaks inversion $\mathcal{P}$ and, in the magnetic system, the $\Theta$ symmetry.
There are four distinct modes transforming under $T_2^-$; these are the axion-active phonons of MBT.

The displacement patterns of the four $T_2^-$ modes involve different subsets of atoms in the septuple layer.
Mode~1 is dominated by Mn displacements, mode~2 involves primarily Bi atoms, mode~3 displaces the inner Te atoms, and mode~4 involves the outer Te atoms.
Because these are zone-boundary modes (at $Z$), the atomic displacements alternate in sign between adjacent unit cells along $c$, ensuring that both inversion and $\Theta$ are broken.

\subsection{Computational details for MBT}

The first-principles calculations for MBT are performed using \code{Quantum ESPRESSO} with fully relativistic PAW pseudopotentials incorporating spin-orbit coupling.
The A-type AFM ordering is imposed via noncollinear magnetism, with the Mn$_1$ and Mn$_2$ sublattices initialized with opposite spin orientations along the $c$ axis.
We employ a plane-wave energy cutoff of 60~Ry (charge density cutoff of 240~Ry) and include van der Waals corrections via the DFT-D3 scheme to properly describe the interlayer bonding.
Calculations are performed with and without a Hubbard $U$ correction on the Mn $3d$ orbitals; we test values of $U = 0$, 1, 3, and 4~eV.

For each axion-active phonon mode, atomic positions are displaced along the normalized \code{SMODES} eigenvector with an amplitude $Q = 0.01$~\AA, chosen to remain within the linear-response regime.
Both the undistorted and displaced structures are converged self-consistently, followed by non-self-consistent (NSCF) calculations on a uniform $k$-point mesh.

The resulting Bloch eigenstates are Wannierized using \code{Wannier90} \cite{marzari1997, marzari2012} with a disentanglement window of $-7$ to $15$~eV and a frozen window of $2.5$ to $11$~eV.
The initial projections consist of Mn $3d$, Bi $6p$, and Te $5p$ orbitals on all atoms in the unit cell, yielding 92 spinor trial states that span 46 Wannier functions (after accounting for spin-orbit coupling) from a set of approximately 75 entangled bands.
The Wannierized tight-binding Hamiltonians are imported into the \code{PythTB} framework \cite{Cole_Python_Tight_Binding_2025} for efficient evaluation of the axion response.

An important technical consideration concerns the $\mathcal{P}\Theta$ symmetry, which enforces Kramers-like degeneracy at every $\mathbf{k}$-point even in the magnetically ordered state.
\code{Quantum ESPRESSO} does not automatically detect this composite symmetry; to enforce the correct degeneracy structure we shift the crystallographic origin to the midpoint between the two Mn sublattices.

\subsection{Results for MBT}

The axion response $\partial_Q \theta$ is computed for each of the four $T_2^-$ modes using \eq{4curv}.
We perform a systematic convergence study as a function of $k$-mesh density, evaluating the integral on uniform meshes ranging from $4 \times 4 \times 4$ to $40 \times 40 \times 40$ points.
The minimum direct band gap of MBT is small (approximately 47~meV at the $Z$ point), requiring dense $k$-meshes for convergence.

For mode~2 of the $T_2^-$ irrep (the Bi-dominated mode) at displacement amplitude $Q = 0.01$~\AA, we obtain a converged value of $\partial_Q \theta \approx 10.0$~rad/\AA\ on the finest meshes.
The convergence behavior is smooth, with the result stabilizing beyond approximately $20 \times 20 \times 20$ $k$-points.

%==========================================================================
\section{Y$_2$Ir$_2$O$_7$}
\seclab{YIO}
%==========================================================================

\subsection{Crystal and magnetic structure}

Y$_2$Ir$_2$O$_7$ (YIO) is a pyrochlore iridate crystallizing in the cubic $Fd\bar{3}m$ space group (No.~227), with conventional lattice parameter $a = 10.19$~\AA\ \cite{wan2011topological, witczak2014correlated}.
The primitive cell contains 22 atoms: 4~Y at the $16d$ Wyckoff position, 4~Ir at $16c$, 2~O at the $8b$ site, and 12~O at the $48f$ site.
The Ir$^{4+}$ ions ($5d^5$, effective total angular momentum $J_\text{eff} = 1/2$) reside on a pyrochlore lattice---a network of corner-sharing tetrahedra---and form corner-sharing IrO$_6$ octahedra with a tilt angle of approximately $53^\circ$.

Below approximately $T_N \approx 150$~K, the Ir moments order in an all-in-all-out (AIAO) magnetic structure, in which the four Ir moments on each tetrahedron point either all toward or all away from the tetrahedral center \cite{wan2011topological, pesin2010mott}.
The AIAO ordering direction on each Ir sublattice points along the local $\langle 111 \rangle$ axis of the pyrochlore lattice.
This magnetic order breaks time-reversal symmetry $\mathcal{T}$ but preserves spatial inversion $\mathcal{P}$, which maps each ``all-in'' tetrahedron to an ``all-out'' tetrahedron and vice versa.
Inversion alone is therefore sufficient to quantize the axion angle to $\theta = 0$ or $\pi$ \cite{Turner2012}.

The phase realized by Y$_2$Ir$_2$O$_7$ with AIAO order depends sensitively on the strength of electronic correlations.
Theoretical calculations predict a Weyl semimetal at weak correlations, a topological axion insulator with $\theta = \pi$ at intermediate correlations, and a trivial Mott insulator with $\theta = 0$ at strong correlations \cite{wan2011topological, witczak2014correlated}.
In this work, we study YIO in a parameter regime (Hubbard $U = 1$--$3$~eV on Ir $5d$) where the system is gapped and in the trivial phase with $\theta = 0$.
This choice provides a non-topological counterpoint to MBT, allowing us to contrast the phonon-induced departure of $\theta$ from $0$ (in YIO) with the departure from $\pi$ (in MBT).
The system exhibits a total magnetization of approximately $2\,\mu_\text{B}$ per formula unit arising from the AIAO-ordered Ir moments.

\subsection{Symmetry analysis and identification of axion-active modes}

In YIO, the axion angle is quantized to $\theta = 0$ solely by spatial inversion $\mathcal{P}$.
An axion-active phonon must therefore break inversion symmetry in order to allow a departure of $\theta$ from zero.
Since $\mathcal{P}$ is a point-group symmetry (with no translational component), it can be broken by zone-center ($\Gamma$-point) phonon modes belonging to irreducible representations that are odd under inversion.

Using the \code{SMODES} software package with the cubic $Fd\bar{3}m$ structure, we enumerate all $\Gamma$-point phonon irreps and their symmetry properties.
The full zone-center phonon representation decomposes as
\begin{equation}
\Gamma_\text{phonon} = A_{1g} \oplus E_g \oplus 2T_{1g} \oplus 4T_{2g} \oplus 3A_{2u} \oplus 3E_u \oplus 8T_{1u} \oplus 4T_{2u},
\eqlab{YIO-decomp}
\end{equation}
where the subscripts $g$ (gerade) and $u$ (ungerade) denote even and odd parity under inversion, respectively.
Only the ungerade representations break $\mathcal{P}$; since the AIAO magnetic order preserves inversion and inversion is the sole protecting symmetry, \emph{any} inversion-odd phonon mode is axion-active.

The ungerade irreps are $A_{2u}$ ($\Gamma_2^-$), $E_u$ ($\Gamma_3^-$), $T_{1u}$ ($\Gamma_4^-$), and $T_{2u}$ ($\Gamma_5^-$).
Among these, only the $T_{1u}$ modes are infrared (IR) active (including the three acoustic translational modes), while the remaining ungerade irreps are optically silent.
The even-parity irreps $A_{1g}$, $E_g$, and $T_{2g}$ are Raman-active; the $T_{1g}$ modes are neither IR- nor Raman-active.

Of particular interest is the $A_{2u}$ ($\Gamma_2^-$) representation, a singly degenerate inversion-odd irrep.
This representation is \emph{neither infrared-active nor Raman-active}: it is a ``silent mode'' that cannot be excited by uniform electric fields or detected by first-order optical spectroscopy.
The \code{SMODES} analysis identifies three distinct $\Gamma_2^-$ modes: one involves displacements of the four Y atoms along their respective local $\langle 111 \rangle$ directions, one displaces the four Ir atoms in an analogous pattern, and one involves the twelve $48f$ oxygen atoms.
In each case, the displacement pattern has the characteristic form of atoms on each sublattice moving along alternating $\langle 111 \rangle$ directions such that inversion symmetry is broken while the cubic point group operations that commute with inversion are preserved.
All three modes are axion-active.

The fact that the $\Gamma_2^-$ mode is optically silent has a significant physical consequence: the axion angle can be modulated by a lattice vibration that is invisible to standard infrared absorption and Raman scattering experiments.
This represents a fundamentally different kind of coupling between the electromagnetic response and lattice degrees of freedom---one that is intrinsically topological in nature and does not manifest in the conventional optical response.

By contrast, the $T_{1u}$ modes are both axion-active and infrared-active, meaning they couple simultaneously to the axion angle and to uniform electric fields.
The $E_u$ and $T_{2u}$ modes are also axion-active but, like $A_{2u}$, are optically silent.
A complete enumeration of the zone-center phonon irreps, their degeneracies, optical activity, and axion activity in YIO is given in Table~\ref{tab:YIO-irreps}.

\begin{table}
\caption{Zone-center phonon irreducible representations of Y$_2$Ir$_2$O$_7$ in the $Fd\bar{3}m$ space group. For each irrep, we list the degeneracy, total number of modes (degeneracy $\times$ multiplicity), IR and Raman activity, and whether the mode is axion-active (i.e., breaks the inversion symmetry that quantizes $\theta = 0$). The $T_{1u}$ count of 24 includes 3 acoustic (translational) modes.}
\label{tab:YIO-irreps}
\begin{tabular}{lcccccc}
\hline\hline
Irrep & Deg. & Mult. & Modes & IR & Raman & Axion \\
\hline
$A_{1g}$ & 1 & 1 & 1 & No & Yes & No \\
$E_{g}$ & 2 & 1 & 2 & No & Yes & No \\
$T_{1g}$ & 3 & 2 & 6 & No & No & No \\
$T_{2g}$ & 3 & 4 & 12 & No & Yes & No \\
$A_{2u}$ & 1 & 3 & 3 & No & No & Yes \\
$E_{u}$ & 2 & 3 & 6 & No & No & Yes \\
$T_{1u}$ & 3 & 8 & 24 & Yes & No & Yes \\
$T_{2u}$ & 3 & 4 & 12 & No & No & Yes \\
\hline\hline
\end{tabular}
\end{table}

\subsection{Computational details for YIO}

The first-principles calculations for Y$_2$Ir$_2$O$_7$ are performed using \code{Quantum ESPRESSO} with fully relativistic PAW pseudopotentials and spin-orbit coupling.
The AIAO magnetic order is imposed by initializing the four Ir sublattice moments along their respective local $\langle 111 \rangle$ directions: Ir$_1$ along $[\bar{1}\bar{1}\bar{1}]$, Ir$_2$ along $[\bar{1}11]$, Ir$_3$ along $[1\bar{1}1]$, and Ir$_4$ along $[11\bar{1}]$.
The four Ir sites are treated as distinct atomic types in the \code{Quantum ESPRESSO} input to allow independent specification of the magnetic moment directions.
We use a plane-wave energy cutoff of 50~Ry (charge density cutoff of 500~Ry) and a $4 \times 4 \times 4$ $k$-point mesh for the self-consistent calculation.
A Hubbard $U$ correction is applied to the Ir $5d$ states within the ortho-atomic DFT+$U$ framework; we test values of $U = 1.0$ and $3.0$~eV.
Marzari-Vanderbilt cold smearing with a broadening of 5~mRy is used to stabilize the self-consistent cycle.

For each $\Gamma_2^-$ phonon mode, we displace the atomic positions by an amplitude of $Q = 0.01$~\AA\ along the normalized \code{SMODES} eigenvector.
The electronic structure of both undistorted and displaced configurations is then Wannierized and the axion response computed following the same four-curvature procedure used for MBT.

%==========================================================================
\section{Computational methodology}
\seclab{methods}
%==========================================================================

\subsection{Phonon symmetry analysis}

The identification of axion-active phonon modes is performed using the \code{SMODES} software package from the ISOTROPY suite \cite{smodes}, which enumerates the irreducible representations of phonon displacement patterns at specified wavevectors for a given space group and Wyckoff positions.
For MBT, we use the non-magnetic hexagonal $R\bar{3}m$ structure and analyze modes at both the $\Gamma$ and $Z$ points.
For YIO, we use the cubic $Fd\bar{3}m$ structure and analyze modes at the $\Gamma$ point.
In each case, the output provides normalized displacement vectors for each mode within each irrep, which serve as the displacement patterns for the first-principles calculations.

\subsection{Density-functional theory calculations}

All density-functional theory calculations are performed using the \code{Quantum ESPRESSO} package with fully relativistic projector augmented-wave (PAW) pseudopotentials, noncollinear magnetism, and spin-orbit coupling.
The exchange-correlation functional is treated within the PBE generalized gradient approximation.
A Hubbard $U$ correction is applied to the transition-metal $d$ states to account for on-site Coulomb repulsion; the optimal value of $U$ is determined by comparison with experimental band gaps and magnetic moments.

For each material, we perform self-consistent (SCF) calculations for the undistorted structure and for structures displaced along each axion-active phonon mode.
Non-self-consistent (NSCF) calculations are then performed on a uniform $k$-point mesh to obtain the Bloch eigenstates needed for Wannierization.

\subsection{Wannierization}

Starting from the NSCF Bloch eigenstates, we construct maximally localized Wannier functions using \code{Wannier90} \cite{marzari1997, marzari2012, soluyanov2011}.
The disentanglement procedure is employed to extract a set of smoothly connected bands from the entangled manifold.
Initial projections are chosen based on the atomic orbital character of the relevant bands (transition-metal $d$ states, heavy-element $p$ states, etc.), and the resulting tight-binding Hamiltonians faithfully reproduce the low-energy electronic structure.

For MBT, the Wannier basis consists of 92 spinor trial states (Mn $3d$, Bi $6p$, Te $5p$) spanning 46 Wannier functions with a disentanglement window of $[-7, 15]$~eV and a frozen window of $[2.5, 11]$~eV.
The Wannierized Hamiltonians are imported into the \code{PythTB} framework \cite{Cole_Python_Tight_Binding_2025}.

\subsection{Evaluation of the axion response}

The axion response $\partial_\lambda \theta$ is evaluated using \eq{4curv} on the Wannierized tight-binding Hamiltonians.
The four-dimensional Berry curvature $\Omega_{ij}$ in the extended space $(k_x, k_y, k_z, \lambda)$ is constructed as follows.

Derivatives of the Bloch states with respect to crystal momenta $k_\mu$ ($\mu = x,y,z$) are obtained from the velocity operators,
\begin{equation}
v_\mu = \frac{1}{\hbar} \partial_{k_\mu} H(\mathbf{k}),
\eqlab{velocity}
\end{equation}
which are computed analytically from the tight-binding Hamiltonian $H(\mathbf{k})$.
Derivatives with respect to the phonon amplitude $\lambda$ are obtained by finite difference,
\begin{equation}
\partial_\lambda | u_{n\mathbf{k}} \rangle \approx \frac{| u_{n\mathbf{k}}(\lambda + \delta\lambda) \rangle - | u_{n\mathbf{k}}(\lambda - \delta\lambda) \rangle}{2\delta\lambda},
\eqlab{findiff}
\end{equation}
where $\delta\lambda$ corresponds to the phonon displacement amplitude $Q$.
The Berry connections and curvatures in the mixed $(k, \lambda)$ space are then assembled from these ingredients, and the Levi-Civita contraction in \eq{4curv} is evaluated numerically on a uniform 3D $k$-mesh.

An important feature of this approach is that it requires only tight-binding Hamiltonians at two (or more) values of $\lambda$, corresponding to the undistorted and displaced crystal structures.
No perturbation-theory machinery within the DFT code is needed; all quantities are computed from the ground-state electronic structures at each displacement.

\subsection{Convergence considerations}

The convergence of $\partial_\lambda \theta$ with respect to $k$-mesh density is controlled by the minimum band gap of the system: smaller gaps require denser meshes, as the Berry curvature is concentrated near gap-closing points.
For MBT, the minimum direct gap is approximately 47~meV at the $Z$ point, necessitating $k$-meshes of at least $20 \times 20 \times 20$ for reliable convergence.
We perform systematic convergence studies with mesh sizes from $4^3$ to $40^3$ for each mode.

%==========================================================================
\section{Summary and outlook}
\seclab{conclusions}
%==========================================================================

We have developed a first-principles methodology for identifying and computing the coupling of the Chern-Simons axion angle to lattice vibrations in magnetic insulators.
By combining symmetry analysis with a gauge-invariant four-curvature formulation evaluated on Wannierized tight-binding Hamiltonians, we avoid the gauge ambiguities inherent in the Chern-Simons 3-form and obtain numerically stable results for $\partial_\lambda \theta$.

We have applied this approach to two materials with contrasting axion topology.
In MnBi$_2$Te$_4$, where $\theta = \pi$ is protected by both inversion $\mathcal{P}$ and the combined symmetry $\Theta = \mathcal{T} S_{1/2}$, the axion-active phonons are zone-boundary modes at the $Z$ point belonging to the $T_2^-$ irreducible representation.
These modes must reside at the zone boundary because the $\Theta$ symmetry involves a half-lattice translation, which can only be broken by distortions that modulate at the corresponding wavevector.
For the Bi-dominated $T_2^-$ mode, we find a substantial axion-phonon coupling of $\partial_Q \theta \approx 10$~rad/\AA, representing a departure from $\pi$.

In Y$_2$Ir$_2$O$_7$, studied in a non-topological regime where $\theta = 0$ is quantized solely by inversion in the AIAO magnetic structure, the axion-active phonons are zone-center modes belonging to any inversion-odd irreducible representation.
We have focused on the $\Gamma_2^-$ ($A_{2u}$) representation, a singly degenerate mode that is notable for being \emph{neither infrared-active nor Raman-active}---a ``silent'' modulator of the axion angle.
This illustrates that the axion-phonon coupling can activate degrees of freedom that are invisible to conventional optical probes.

The two materials studied here illustrate the diversity of axion-phonon physics across both topological and trivial magnetic insulators.
In MBT, the axion-active modes are necessarily at the zone boundary due to the translational component of the protecting symmetry and modulate $\theta$ away from $\pi$, while in YIO they are at the zone center, optically silent, and modulate $\theta$ away from $0$.
Comparing the magnitudes and character of the axion response in these two regimes provides insight into how the axion-phonon coupling depends on the underlying band topology.
These results motivate further exploration of axion-phonon coupling across different magnetic insulators, as well as investigations into whether such couplings could be detected through alternative experimental probes such as neutron scattering, hyper-Raman spectroscopy, or nonlinear optical measurements.

\bibliography{ref}

\end{document}
